\section{Conclusion}
\label{sec:conclusion}

We reframed Quiz \#3 as a structured measurement exercise aligned with the NIST AI RMF trustworthy characteristics \citep{nist2023airmf}: each quiz item became a research question with a hypothesis, a logged single-turn experiment, and a rubric label. Under the illustrative score template bundled in \texttt{output/results/nist\_quiz\_scores.json}, the model appears strongest on explicit safety refusals, prompt-injection resistance, and several accountability prompts, while validity (base rates, citations) and operational transparency items cluster in partial compliance.

\textbf{Next steps.} Run \texttt{python scripts/run\_nist\_llm\_evaluation.py} with your API credentials, adjudicate \{C,P,N\} honestly, refresh the JSON score file, regenerate \texttt{nist\_rubric\_table.tex}, and recompile. Repeating the battery across temperatures or model versions would support a measurement-style discussion of variance, which a single snapshot cannot provide.
